\section{Conclusion}

The preceding discussion interpreted the findings against the theoretical framework of Section 2 and the five hypotheses. This conclusion ties the findings together, states the boundaries of the replication claim, and looks ahead to future research. This study set out to replicate Beirne and Sugandi (2023) for Japan and to test whether the paper's findings survive an independent rebuild and an extended sample. The replication estimated the paper's recursively restricted structural VAR on working-days data from 14 January 1999 to 31 March 2021 with 2021-era data vintages, added the unrestricted specification as a robustness check, and extended both through 30 June 2026 with current-vintage data. The paper's core results reproduce. The yen appreciates in real effective terms after a risk-off shock, the JGB spread widens relative to the United States, real GDP declines with a lag, and portfolio flows do not respond significantly. Nine of the seventeen variable-frequency comparisons against the published figures are matches, five are partial matches, and three are mismatches, and the mismatches have documented mechanisms.

The three research questions posed in the introduction are answered as follows. The paper's qualitative structure reproduces over the original sample, with the RGDP, REER, spread, and capital flow results all confirmed and the WUI and other investment responses failing to reproduce. The relationships survive into the post-2021 regime only partially, with the price channels persisting and the real channel attenuating to statistical indistinguishability from zero at the monthly frequency. The results are vintage-dependent in exactly one place that matters, the RGDP response, which flips sign on current-vintage data, and the dual-vintage policy is what restores the paper's finding.

The five hypotheses of Section 2 resolve into three rejections and two non-rejections. The nulls of no significant REER response, no significant RGDP response in the paper period, and no significant spread response are rejected. The nulls of no significant portfolio flow response and no significant uncertainty response are not rejected, matching the paper's conclusions for those channels, with the direct investment response at the three-month horizon as the one qualified exception.

The objectives set out in the introduction were met. The paper-period specification was rebuilt and validated against the published figures, the extended sample was estimated and documented, and every specification choice, code defect, and data-vintage decision was disclosed in the text. The three contributions stand. The validation confirms the paper's core results for Japan under an independent implementation. The extension shows that the output channel weakened after 2021 in a way that the post-2021 literature would predict. And the documentation of the five corrected code defects and the vintage-induced sign reversal gives the next replication a map of the traps.

What the study does not claim is equally explicit. The Swiss and United States models are not re-estimated, the weekly frequency panels are not reproduced, and the uncertainty and other investment responses could not be matched to the published figures. Each omission is either a scope decision stated in the Research Design or a documented failure listed in the limitations, and each appears again in the future research directions. The replication claim is bounded by these boundaries, and the reader is told exactly where they sit.

\subsection{Recommendations for Future Research}

The boundaries of this replication define the agenda for the work that follows it. What the study did not do is stated directly, and each omission maps to a concrete direction that others can undertake.

\textit{What this study did not do.}

\begin{itemize}
  \item The weekly frequency panels of the paper's Appendix 4 were not replicated. The two-tier design covers the daily and monthly frequencies only.
  \item The Swiss and United States models were not estimated. The paper's cross-country comparison, and its claim that the negative real spillover is unique to Japan, remain untested by an independent rebuild.
  \item The uncertainty response could not be reproduced at either frequency. The daily and monthly WUI comparisons both fail.
  \item The daily other investment response could not be reproduced. The sign and phase of the response differ from the paper's figure.
  \item Exact magnitudes could not be reproduced. The spread runs about 1.5 times the paper's peak and the REER plateau 2.5 times, because the paper's EViews settings are unreported.
\end{itemize}

\textit{What others can do.}

\begin{itemize}
  \item Estimate the weekly tier to complete the frequency comparison and test whether the extended-period attenuation holds at the intermediate frequency.
  \item Estimate the full three-country system to test whether the paper's cross-country asymmetries survive an independent implementation.
  \item Resolve the uncertainty mismatch with a higher-frequency uncertainty measure, or with a construction that separates global from foreign-specific episodes, since the sign of the response depends on the episode type.
  \item Compare the Bank of Japan BPM6 and the IMF IFS capital flow sources directly on the same sample to settle whether the other investment phase reversal is a source-data effect.
  \item Investigate the direct investment response at the three-month horizon, the only capital flow finding that contradicts the paper, and the boundary between short-term portfolio dynamics and long-horizon investment decisions.
  \item Model the conditional safe-haven behavior after 2021 with regime-switching or interaction specifications that let the impulse responses vary with the US-Japan rate differential.
  \item Run an event study on the August 2024 carry trade unwind with high-frequency data that the balance of payments cannot capture.
\end{itemize}

Each direction is a stopping point that the present data and methods cannot reach, and each marks the boundary of what this replication can claim.
